\documentclass{beamer}
\usepackage{graphicx}
\graphicspath{ {./images/} }
% Theme choice:
\usetheme{Antibes}



\title{TAMAGOTCHI GAME USING LUA} 
\author{Tamagoshe - Sharon, Ananya, Kavya, Shreya}
\logo{\large \LaTeX{}}

\begin{document}

\begin{frame}
    \titlepage 
\end{frame}

% Remove logo from the next slides
\logo{}

\section{MODIFICATIONS THAT WE ARE MAKING}
\begin{frame}{Device Shape}
\begin{itemize}
    \item  A duck-shaped handheld device.
    \item User Input: Buttons A, B, C
\end{itemize}
\begin{figure}[h]
    \includegraphics[width=0.4\textwidth]{pic/tamago.png} % adjust width as needed
    \label{fig:tamago}
\end{figure}
\end{frame}

\begin{frame}{Pet}
\begin{itemize}
    \item NAME: WEP
    \item DESCRIPTION: A digital character who is
    passionate about coding, software development, and technology. WEP has a quirky personality and loves solving puzzles and tackling challenges.
    \item FOOD:
        \begin{enumerate}
            \item Biscuits: Keeps WEP satisfied and happy. 
            \item Chai: Provides energy and boosts creativity, essential for coding sessions.
        \end{enumerate}
    \item MEDICINE: Pills: Helps WEP recover when she's feeling unwell.
    \item CLEAN: Deodorant: Keeps WEP fresh and clean. 
\end{itemize}
\end{frame}

\begin{frame}{Games}
\begin{enumerate}
    \item Python Perfection: WEP identifies and fixes code errors and guess outputs
    \item Hackathon Prep: WEP has to identify the tech stacks + format problem statements, etc.
    \item Catch the Duck:  WEP chases a duck, providing a fun break from coding.
\end{enumerate}
\end{frame}

\begin{frame}{Meters}
\begin{itemize}
    \item Coder Meter:
    \begin{itemize}
        \item Description: Reflects WEP's coding proficiency and engagement.
        \item Increase by: Playing coding-related games like Guess the Code Error and Hackathon Prep.
        \item Decrease by: Lack of coding activities.
    \end{itemize}
    \item Happiness Meter:
    \begin{itemize}
        \item Description: Shows how happy and content WEP is.
        \item Increase by: Playing any game, eating biscuits, and drinking chai.
        \item Decrease by: Neglect, lack of playtime, and lack of food.
    \end{itemize}
\end{itemize}
\end{frame}

\begin{frame}{Meters}
\begin{itemize}
    \item Health Meter:
    \begin{itemize}
        \item Description: Indicates WEP's overall health and well-being.
        \item Increase by: Feeding chai and biscuits, cleaning with deodorant, and ensuring proper sleep by turning off the light.
        \item Decrease by: Poor diet, lack of cleanliness, and insufficient sleep.
    \end{itemize}
\end{itemize}
\end{frame}

\begin{frame}{Game Timer}
\begin{itemize}
    \item Game Timer:  
    \begin{itemize}
        \item The game operates on a 15-minute cycle, during which the three meters gradually decrease. Players need to regularly interact with WEP to keep her engaged, happy, and healthy.
    \end{itemize}
\end{itemize}
\end{frame}

\end{document}
